\section{DIAL: ampliar la cobertura lingüística de un Offline Dataset con VLM}

RT-1 demostró que la imitación fuera de línea a gran escala es viable, pero dejó al descubierto un problema de economía de datos: 130,000 episodes, 13 robots, 17 meses y unas 700 tasks, además de la necesidad de reescribir en lenguaje natural los teleoperator command estructurados. DIAL no vuelve a recolectar acciones, sino que usa un VLM para generar instruction labels más ricas para las trayectorias existentes.

\subsection{Agenda y Teleoperation Cost}

Los resultados de scaling de RT-1 llevan con facilidad a concentrar la atención en la architecture y a pasar por alto el costo de recolección y de anotación lingüística que hay detrás de cada robot trajectory. El punto de partida de DIAL es precisamente volver a calcular el costo de esta cadena de suministro de datos: qué pasos deben realizar el robot y el operador, y qué trabajo semántico puede delegarse a un modelo de visión y lenguaje preentrenado. Esta sección explica primero los dos niveles de costo humano del pipeline original; estas cifras explican el cuello de botella de la escalabilidad, pero de ellas no se puede concluir directamente que las etiquetas automáticas tengan la misma precisión que las etiquetas humanas.

La Agenda entra en DIAL. La slide 37 divide la cadena de datos original en dos etapas: el operador recolecta acciones con comandos estructurados y después un crowd worker reescribe los comandos en lenguaje natural. Los dos niveles de costo humano hacen que cada semantic concept adicional resulte muy caro y que la gran mayoría de las trayectorias solo tenga etiquetas basadas en plantillas.

\lecturefigure{slide-36-agenda-dial.jpg}{Agenda: DIAL entra en la capa de data augmentation}{Diapositiva V036 recuperada de la grabación oficial; Stanford CS25 Lecture 15}

\lecturefigure{slide-37-rt1-teleoperation-cost.jpg}{Economía de datos de RT-1: 130k episodes, 13 robots, 17 months, 700 tasks y anotación lingüística humana}{Diapositiva V037 recuperada de la grabación oficial; Stanford CS25 Lecture 15}

\teachervoice{El ponente plantea la motivación de DIAL de forma muy pragmática: el «secret ingredient» de RT-1 no es alguna fancy layer, sino una ingeniería de datos costosa y prolongada. Si solo se habla del modelo y no de quién hace la teleoperación ni de quién reetiqueta, no es posible explicar el costo real del scaling en robótica.}

\subsection{Pipeline de cuatro etapas}

DIAL primero hace fine-tune de un VLM con un conjunto pequeño de crowd-sourced instructions; después lo usa para generar lenguaje para un offline dataset de gran escala; a continuación mezcla las trayectorias originales, las etiquetas humanas y las etiquetas del VLM para entrenar una language-conditioned policy; por último, evalúa un RT-1-style controller con lenguaje natural que no apareció durante el entrenamiento.

\lecturefigure{slide-38-dial-pipeline.jpg}{Las cuatro etapas de DIAL: poco lenguaje humano, VLM relabeling, entrenamiento mixto y unseen instruction evaluation}{Diapositiva V038 recuperada de la grabación oficial; Xiao et al., 2022}

Sea $\mathcal{D}_A$ el conjunto pequeño de etiquetas humanas en lenguaje natural, $\mathcal{D}_B$ el conjunto grande de trayectorias estructuradas originales y $\hat{\mathcal{D}}_B$ las trayectorias reetiquetadas por el VLM; entonces el objetivo de la policy puede escribirse como
\begin{equation}
\mathcal{L}_{\mathrm{DIAL}}(\theta)
=-\mathbb{E}_{(o,l,a)\sim
\mathcal{D}_A\cup\mathcal{D}_B\cup\hat{\mathcal{D}}_B}
\log p_{\theta}(a\mid o,l).
\end{equation}
Lo importante no es la complejidad de la fórmula, sino que se amplía el soporte de $l$: un mismo comportamiento puede recibir relaciones espaciales, atributos, expresiones sinónimas y descripciones más naturales.

\subsection{El Language Space pasa de Cluster a cobertura}

Los embedding clusters de la izquierda provienen de los comandos estructurados; los colores y los límites de los clústeres son muy nítidos, porque cada plantilla equivale aproximadamente a un one-hot task ID. Las DIAL predictions de la derecha son más continuas, y las distintas expresiones se entremezclan en el espacio semántico. Los ejemplos del centro muestran cómo trayectorias como `open top drawer` o `move bottle near chip bag` se reescriben con un lenguaje más descriptivo.

\lecturefigure{slide-39-dial-language-clusters.jpg}{Comparación en el embedding-space entre Structured commands y DIAL predicted language}{Diapositiva V039 recuperada de la grabación oficial; Xiao et al., 2022}

\begin{knowledgebox}{Para qué sirve un Language Space denso}
Las etiquetas de plantilla solo le indican a la policy el «número de tarea», y difícilmente permiten compartir conceptos como `left`, `lonely object`, colores o formas. Un lenguaje más rico permite que distintas tareas reutilicen representaciones en el nivel del significado de las palabras, y que en la prueba se especifique la misma acción con una paraphrase no vista en el entrenamiento. El costo es que el VLM produce ruido e incluso descripciones erróneas.
\end{knowledgebox}

\subsection{Novel Instructions: ver el total y también las categorías}

La slide de resultados compara a la izquierda distintas dataset mixtures y muestra a la derecha instrucciones no vistas como `raise the leftmost can`, `move the lonely object to the others` y `move the right apple to the left`. El DIAL completo alcanza 60\% en el overall; las categories spatial y rephrased son las más sólidas, mientras que la semantic category sigue siendo claramente más difícil.

\lecturefigure{slide-40-dial-novel-instruction-results.jpg}{DIAL novel-instruction evaluation: combinaciones de datos, resultados por categoría y ejemplos cualitativos}{Diapositiva V040 recuperada de la grabación oficial; Xiao et al., 2022}

\begin{warningbox}{60\% no equivale a control con lenguaje abierto}
La prueba incluye poco más de sesenta novel instructions construidas a mano, que cubren variaciones espaciales, reformulaciones y algunos cambios semánticos, pero no lenguaje natural arbitrario. El ponente también afirma explícitamente que la ablation sistemática de la compositional generalization está incompleta y que el mecanismo interno de muchos casos exitosos no está claro.
\end{warningbox}

\subsection{Takeaway: con suficiente Diversity, se tolera poco Label Noise}

Los DIAL takeaways son tres: usar el VLM como data augmentation; inyectar Internet-scale priors en los offline data; y, cuando el dataset es lo bastante grande y los comportamientos son diversos, un poco de ruido en las etiquetas puede seguir siendo aceptable. Este último punto no significa que la label quality sea irrelevante, sino que las trayectorias redundantes y la diversidad semántica pueden amortiguar parte de los errores.

\lecturefigure{slide-41-dial-takeaways.jpg}{DIAL takeaways: VLM augmentation, Internet priors, offline diversity y label noise limitado}{Diapositiva V041 recuperada de la grabación oficial; Stanford CS25 Lecture 15}

\teachervoice{Ted Xiao es muy cauteloso con «label noise okay»: un ruido demasiado alto sigue dañando el entrenamiento, y el propio comportamiento del robot debe ser lo bastante diverso. Ni el mejor captioner puede crear de la nada una habilidad no recolectada a partir de trayectorias de acción completamente idénticas.}

\subsection{Resumen de la sección}

DIAL coloca el foundation model en la capa de datos y no en la capa de control: el VLM reescribe la interfaz semántica de las trayectorias existentes, de modo que los mismos datos de acción admiten más lenguaje. Reduce el costo de la anotación humana y mejora el desempeño con novel instruction, pero no puede sustituir la cobertura de comportamientos, y la evidencia sobre la generalización composicional sigue siendo limitada.
